\documentclass[../../../include/open-logic-section]{subfiles}

\begin{document}

\olfileid{his}{set}{cantorplane}
\olsection{د خط او مستوي په اړه د کانتور کار}

د شروډر--برنشتاين د ثبوت ځينې اړوند حالات له هغه تاريخ سره تړاو لري چې په \olref[his][set][pathology]{sec} کښې مو بيان کړ۔ په ياد ولرئ چې په 1877 کښې کانتور ثابته کړه چې په يوه مربع کښې عين يو هومره نقطې دي څومره چې يې په يوې څنډې کښې دي۔ دلته به د هغه ثبوت (لومړۍ هڅه) وړاندې کړو۔

فرض کړئ چې $\unitline$ واحد خط وي، يعنې د نقطو سټ $[0,1]$۔ فرض کړئ چې $\unitsquare$ واحده مربع وي، يعنې د نقطو سټ $\unitline
\times \unitline$۔ په دې ټکو، کانتور ثابته کړه چې $\cardeq{\unitline}{\unitsquare}$۔ هغه ډېډېکنډ ته يو ليک وليکه چې په اصل کښې يې لاندې استدلال درلود۔

\begin{thm}\ollabel{thm:cantorplane}
$\cardeq{\unitline}{\unitsquare}$
\end{thm}

\begin{proof}[ثبوت: لومړۍ برخه۔]
$a, b \in \unitline$ ثابت کړئ۔ په دوه‌ييزه نښه‌ليکنه يې وليکئ، څو د $0$ونو او $1$ونو بې‌پايه لړۍ، $a_1$، $a_2$، \dots، او $b_1$، $b_2$، \dots، ولرو، داسې چې:
\begin{align*}
a &= 0.a_1a_2a_3a_4\dots\\
b &= 0.b_1b_2b_3b_4\dots
\intertext{اوس تابع $f \colon \unitsquare \to \unitline$ په پام کښې ونيسئ، چې داسې راکول کېږي:} 
f(a, b) & = 0.a_1b_1a_2b_2a_3b_3a_4b_4\dots
\end{align*}
اوس $f$ !!a{injection} دے، ځکه که $f(a, b) = f(c,d)$ وي، نو $a_n =
c_n$ او $b_n = d_n$ د هر $n \in \Nat$ لپاره دي؛ ځکه نو $a = c$ او $b =
d$ دي۔
\end{proof}

له بده مرغه، لکه چې ډېډېکنډ کانتور ته څرګنده کړه، دا اصلي پوښتنه نۀ ځوابوي۔ $0.\dot{1}\dot{0} =
0.1010101010\ldots$ په پام کښې ونيسئ۔ اړتيا لرو چې $f(a,b) = 0.\dot{1}\dot{0}$ وي، چې:
\begin{align*}
a&= 0.\dot{1}\dot{1} = 0.111111\ldots\\
b&= 0
\end{align*}
خو $a = 0.\dot{1}\dot{1} = 1$۔ نو چې وايو «$a$ او $b$ په دوه‌ييزه نښه‌ليکنه وليکئ»، بايد وټاکو چې \emph{کومه} نښه‌ليکنه کاروو؛ او ځکه چې $f$ بايد يوه \emph{تابع} وي، د دوو ممکنو نښه‌ليکنو څخه يوازې \emph{يوه} کارولے شو۔ خو که، مثلاً، ساده نښه‌ليکنه وکاروو او $a$ داسې وليکو: «$1.000\ldots$»، نو داسې هېڅ جوړه $\tuple{a, b}$ نۀ لرو چې $f(a, b) = 0.\dot{1}\dot{0}$ وي۔ \emph{سرچينه‌يي يادونه OLSTH-050: د پورتني تصوير طرحه هر عدد په صفر پيلېدونکې کسري بڼه ليکي۔ د وقفې ښے پای، يعنې واحد عدد، په همدغه بڼه د يوونو بې‌پايه تکرار غواړي۔ چاپي ساده بڼه له اعشاري ټکي مخکې واحد لري؛ د هغې د صحيحې برخې غورځول به واحد له صفر سره يو شان کوډ کړي او يو پر يو والے به مات کړي۔ د ټولو نورو دوه‌ګونې بڼې لرونکو عددونو لپاره يو ثابت انتخاب، او د ښي پای لپاره دا ځانګړے تکراري انتخاب، اړين دي۔ اصلي بېلګه ساتل شوې او د پای‌ټکي پاتې قيد څرګند شو۔}

لنډه دا چې ډېډېکنډ څرګنده کړه چې، د ځينو تکراري اعشاري پراختياوو د ممکنوالي له امله، د کانتور تابع $f$ !!a{injection} دے، خو !!a{surjection} \emph{نۀ} دے۔ نو کانتور يوازې $\cardle{\unitsquare}{\unitline}$ ښودلې ده، \emph{نۀ} دا چې $\cardeq{\unitsquare}{\unitline}$۔

کانتور نژدې سمدستي ډېډېکنډ ته ځواب وليکه او په اصل کښې يې وړانديز وکړ چې ثبوت په لاندې ډول بشپړېدای شي:

\begin{proof}[ثبوت: بشپړه برخه۔]
نو موږ ښودلې ده چې $\cardle{\unitsquare}{\unitline}$۔ خو ښکاره ده چې له $\unitline$ څخه $\unitsquare$ ته !!a{injection} شته: يوازې خط د مربع د يوې څنډې پر اوږدوالي کېږدئ۔ نو $\cardle{\unitline}{\unitsquare}$ او $\cardle{\unitsquare}{\unitline}$ دي۔ د شروډر--برنشتاين له مخې (\olref[sfr][siz][sb]{thm:schroder-bernstein})، $\cardeq{\unitline}{\unitsquare}$۔
\end{proof}

خو، څرګنده ده چې کانتور وروستۍ کرښه په دې ټکو نۀ شوه بشپړولے، ځکه د شروډر--برنشتاين قضيه لا نۀ وه ثابته شوې۔ په رښتيا، که څه هم کانتور وروسته دا د يوې عمومي فرضيې په توګه وړاندې کړه، خو تر 1897 پورې يې د قناعت وړ ثبوت نۀ و ورکړل شوے۔ (ځکه نو، وروسته په 1877 کښې، کانتور د \olref{thm:cantorplane} يو بل ثبوت وړاندې کړ چې د شروډر--برنشتاين له لارې نۀ و۔)

\end{document}
